\section{Coherent sheaves on Proj, I}
\label{section-coherent-proj}

\noindent
In this section we discuss coherent sheaves on $\text{Proj}(A)$
where $A$ is a Noetherian graded ring generated by $A_1$ over $A_0$.
In the next section we discuss what happens if $A$ is not generated
by degree $1$ elements. First, we formulate an all-in-one result for
projective space over a Noetherian ring.

\begin{lemma}
\label{lemma-coherent-projective}
Let $R$ be a Noetherian ring.
Let $n \geq 0$ be an integer.
For every coherent sheaf $\mathcal{F}$ on $\mathbf{P}^n_R$
we have the following:
\begin{enumerate}
\item There exists an $r \geq 0$ and
$d_1, \ldots, d_r \in \mathbf{Z}$ and a surjection
$$
\bigoplus\nolimits_{j = 1, \ldots, r}
\mathcal{O}_{\mathbf{P}^n_R}(d_j)
\longrightarrow
\mathcal{F}.
$$
\item We have $H^i(\mathbf{P}^n_R, \mathcal{F}) = 0$ unless
$0 \leq i \leq n$.
\item For any $i$ the cohomology group $H^i(\mathbf{P}^n_R, \mathcal{F})$
is a finite $R$-module.
\item If $i > 0$, then
$H^i(\mathbf{P}^n_R, \mathcal{F}(d)) = 0$ for all $d$ large enough.
\item For any $k \in \mathbf{Z}$ the graded $R[T_0, \ldots, T_n]$-module
$$
\bigoplus\nolimits_{d \geq k} H^0(\mathbf{P}^n_R, \mathcal{F}(d))
$$
is a finite $R[T_0, \ldots, T_n]$-module.
\end{enumerate}
\end{lemma}

\begin{proof}
We will use that $\mathcal{O}_{\mathbf{P}^n_R}(1)$ is an ample invertible
sheaf on
the scheme $\mathbf{P}^n_R$. This follows directly from the definition
since $\mathbf{P}^n_R$ covered by the standard affine opens $D_{+}(T_i)$.
Hence by
Properties, Proposition \ref{properties-proposition-characterize-ample}
every finite type quasi-coherent $\mathcal{O}_{\mathbf{P}^n_R}$-module
is a quotient of a finite direct sum of tensor powers of
$\mathcal{O}_{\mathbf{P}^n_R}(1)$. On the other hand coherent sheaves
and finite type quasi-coherent sheaves are the same thing on projective
space over $R$ by Lemma \ref{lemma-coherent-Noetherian}. Thus we see (1).

\medskip\noindent
Projective $n$-space $\mathbf{P}^n_R$ is covered by $n + 1$ affines,
namely the standard opens $D_{+}(T_i)$, $i = 0, \ldots, n$, see Constructions,
Lemma \ref{constructions-lemma-standard-covering-projective-space}.
Hence we see that for any quasi-coherent
sheaf $\mathcal{F}$ on $\mathbf{P}^n_R$
we have $H^i(\mathbf{P}^n_R, \mathcal{F}) = 0$ for $i \geq n + 1$,
see Lemma \ref{lemma-vanishing-nr-affines}. Hence (2) holds.

\medskip\noindent
Let us prove (3) and (4) simultaneously for all coherent sheaves
on $\mathbf{P}^n_R$ by descending induction on $i$. Clearly the result
holds for $i \geq n + 1$ by (2). Suppose we know the result for
$i + 1$ and we want to show the result for $i$. (If $i = 0$, then
part (4) is vacuous.) Let $\mathcal{F}$ be a coherent sheaf on
$\mathbf{P}^n_R$. Choose a surjection as in (1) and denote
$\mathcal{G}$ the kernel so that we have a short exact sequence
$$
0 \to \mathcal{G} \to
\bigoplus\nolimits_{j = 1, \ldots, r}
\mathcal{O}_{\mathbf{P}^n_R}(d_j)
\to
\mathcal{F} \to 0
$$
By Lemma \ref{lemma-coherent-abelian-Noetherian}
we see that $\mathcal{G}$ is coherent. The long exact
cohomology sequence gives an exact sequence
$$
H^i(\mathbf{P}^n_R, \bigoplus\nolimits_{j = 1, \ldots, r}
\mathcal{O}_{\mathbf{P}^n_R}(d_j))
\to
H^i(\mathbf{P}^n_R, \mathcal{F})
\to
H^{i + 1}(\mathbf{P}^n_R, \mathcal{G}).
$$
By induction assumption the right $R$-module is finite and by
Lemma \ref{lemma-cohomology-projective-space-over-ring} the left
$R$-module is finite. Since $R$ is Noetherian it follows immediately
that $H^i(\mathbf{P}^n_R, \mathcal{F})$ is a finite $R$-module.
This proves the induction step for assertion (3).
Since $\mathcal{O}_{\mathbf{P}^n_R}(d)$ is invertible
we see that twisting on $\mathbf{P}^n_R$ is an exact functor (since
you get it by tensoring with an invertible sheaf, see
Constructions, Definition \ref{constructions-definition-twist}).
This means that for all $d \in \mathbf{Z}$ the sequence
$$
0 \to \mathcal{G}(d) \to
\bigoplus\nolimits_{j = 1, \ldots, r}
\mathcal{O}_{\mathbf{P}^n_R}(d_j + d)
\to
\mathcal{F}(d) \to 0
$$
is short exact. The resulting cohomology sequence is
$$
H^i(\mathbf{P}^n_R, \bigoplus\nolimits_{j = 1, \ldots, r}
\mathcal{O}_{\mathbf{P}^n_R}(d_j + d))
\to
H^i(\mathbf{P}^n_R, \mathcal{F}(d))
\to
H^{i + 1}(\mathbf{P}^n_R, \mathcal{G}(d)).
$$
By induction assumption we see the module on the right is zero
for $d \gg 0$ and by the computation in
Lemma \ref{lemma-cohomology-projective-space-over-ring}
the module on the left is zero as soon as $d \geq -\min\{d_j\}$
and $i \geq 1$. Hence the induction step for assertion (4).
This concludes the proof of (3) and (4).

\medskip\noindent
In order to prove (5) note that for all sufficiently large $d$
the map
$$
H^0(\mathbf{P}^n_R, \bigoplus\nolimits_{j = 1, \ldots, r}
\mathcal{O}_{\mathbf{P}^n_R}(d_j + d))
\to
H^0(\mathbf{P}^n_R, \mathcal{F}(d))
$$
is surjective by the vanishing of $H^1(\mathbf{P}^n_R, \mathcal{G}(d))$
we just proved. In other words, the module
$$
M_k
=
\bigoplus\nolimits_{d \geq k} H^0(\mathbf{P}^n_R, \mathcal{F}(d))
$$
is for $k$ large enough a quotient of the corresponding module
$$
N_k
=
\bigoplus\nolimits_{d \geq k} H^0(\mathbf{P}^n_R,
\bigoplus\nolimits_{j = 1, \ldots, r}
\mathcal{O}_{\mathbf{P}^n_R}(d_j + d)
)
$$
When $k$ is sufficiently small (e.g.\ $k < -d_j$ for all $j$) then
$$
N_k = \bigoplus\nolimits_{j = 1, \ldots, r}
R[T_0, \ldots, T_n](d_j)
$$
by our computations in Section \ref{section-cohomology-projective-space}.
In particular it is finitely generated.
Suppose $k \in \mathbf{Z}$ is arbitrary.
Choose $k_{-} \ll k \ll k_{+}$.
Consider the diagram
$$
\xymatrix{
N_{k_{-}} & N_{k_{+}} \ar[d] \ar[l] \\
M_k & M_{k_{+}} \ar[l]
}
$$
where the vertical arrow is the surjective map above and
the horizontal arrows are the obvious inclusion maps.
By what was said above we see that $N_{k_{-}}$ is a finitely
generated $R[T_0, \ldots, T_n]$-module. Hence $N_{k_{+}}$ is
a finitely generated $R[T_0, \ldots, T_n]$-module because it
is a submodule of a finitely generated module and the ring
$R[T_0, \ldots, T_n]$ is Noetherian. Since the vertical arrow
is surjective we conclude that $M_{k_{+}}$ is a finitely
generated $R[T_0, \ldots, T_n]$-module. The quotient
$M_k/M_{k_{+}}$ is finite as an $R$-module since it is a
finite direct sum of the finite $R$-modules
$H^0(\mathbf{P}^n_R, \mathcal{F}(d))$ for $k \leq d < k_{+}$.
Note that we use part (3) for $i = 0$ here. Hence
$M_k/M_{k_{+}}$ is a fortiori a finite $R[T_0, \ldots, T_n]$-module.
In other words, we have sandwiched $M_k$ between two finite
$R[T_0, \ldots, T_n]$-modules and we win.
\end{proof}

\begin{lemma}
\label{lemma-coherent-on-proj}
Let $A$ be a graded ring such that $A_0$ is Noetherian and
$A$ is generated by finitely many elements of $A_1$ over $A_0$.
Set $X = \text{Proj}(A)$. Then $X$ is a Noetherian scheme.
Let $\mathcal{F}$ be a coherent $\mathcal{O}_X$-module.
\begin{enumerate}
\item There exists an $r \geq 0$ and
$d_1, \ldots, d_r \in \mathbf{Z}$ and a surjection
$$
\bigoplus\nolimits_{j = 1, \ldots, r} \mathcal{O}_X(d_j)
\longrightarrow \mathcal{F}.
$$
\item For any $p$ the cohomology group $H^p(X, \mathcal{F})$ is a finite
$A_0$-module.
\item If $p > 0$, then $H^p(X, \mathcal{F}(d)) = 0$ for all $d$ large enough.
\item For any $k \in \mathbf{Z}$ the graded $A$-module
$$
\bigoplus\nolimits_{d \geq k} H^0(X, \mathcal{F}(d))
$$
is a finite $A$-module.
\end{enumerate}
\end{lemma}

\begin{proof}
By assumption there exists a surjection of graded $A_0$-algebras
$$
A_0[T_0, \ldots, T_n] \longrightarrow A
$$
where $\deg(T_j) = 1$ for $j = 0, \ldots, n$. By Constructions, Lemma
\ref{constructions-lemma-surjective-graded-rings-generated-degree-1-map-proj}
this defines a closed immersion $i : X \to \mathbf{P}^n_{A_0}$
such that $i^*\mathcal{O}_{\mathbf{P}^n_{A_0}}(1) = \mathcal{O}_X(1)$.
In particular, $X$ is Noetherian as a closed subscheme of the Noetherian
scheme $\mathbf{P}^n_{A_0}$. We claim that the results of the lemma for
$\mathcal{F}$ follow from the corresponding
results of Lemma \ref{lemma-coherent-projective} for the coherent sheaf
$i_*\mathcal{F}$ (Lemma \ref{lemma-i-star-equivalence}) on
$\mathbf{P}^n_{A_0}$. For example, by this lemma there
exists a surjection
$$
\bigoplus\nolimits_{j = 1, \ldots, r} \mathcal{O}_{\mathbf{P}^n_{A_0}}(d_j)
\longrightarrow i_*\mathcal{F}.
$$
By adjunction this corresponds to a map
$\bigoplus_{j = 1, \ldots, r} \mathcal{O}_X(d_j) \longrightarrow \mathcal{F}$
which is surjective as well. The statements on cohomology follow from the
fact that
$H^p(X, \mathcal{F}(d)) = H^p(\mathbf{P}^n_{A_0}, i_*\mathcal{F}(d))$
by Lemma \ref{lemma-relative-affine-cohomology}.
\end{proof}

\begin{lemma}
\label{lemma-recover-tail-graded-module}
Let $A$ be a graded ring such that $A_0$ is Noetherian and $A$ is generated
by finitely many elements of $A_1$ over $A_0$. Let $M$ be a
finite graded $A$-module. Set $X = \text{Proj}(A)$ and let $\widetilde{M}$
be the quasi-coherent $\mathcal{O}_X$-module on $X$ associated to $M$.
The maps
$$
M_n \longrightarrow \Gamma(X, \widetilde{M}(n))
$$
from Constructions, Lemma \ref{constructions-lemma-apply-modules}
are isomorphisms for all sufficiently large $n$.
\end{lemma}

\begin{proof}
Because $M$ is a finite $A$-module we see that
$\widetilde{M}$ is a finite type $\mathcal{O}_X$-module,
i.e., a coherent $\mathcal{O}_X$-module.
Set $N = \bigoplus_{n \in \mathbf{Z}} \Gamma(X, \widetilde{M}(n))$.
We have to show that the map $M \to N$ of graded $A$-modules
is an isomorphism in all sufficiently large degrees.
By Properties, Lemma \ref{properties-lemma-proj-quasi-coherent}
we have a canonical isomorphism $\widetilde{N} \to \widetilde{M}$
such that the induced maps $N_n \to N_n = \Gamma(X, \widetilde{M}(n))$
are the identity maps. Thus we have maps
$\widetilde{M} \to \widetilde{N} \to \widetilde{M}$
such that for all $n$ the diagram
$$
\xymatrix{
M_n \ar[d] \ar[r] & N_n \ar[d] \ar@{=}[rd] \\
\Gamma(X, \widetilde{M}(n)) \ar[r] &
\Gamma(X, \widetilde{N}(n)) \ar[r]^{\cong} &
\Gamma(X, \widetilde{M}(n))
}
$$
is commutative. This means that the composition
$$
M_n \to \Gamma(X, \widetilde{M}(n))
\to \Gamma(X, \widetilde{N}(n)) \to \Gamma(X, \widetilde{M}(n))
$$
is equal to the canonical map $M_n \to \Gamma(X, \widetilde{M}(n))$.
Clearly this implies that the composition
$\widetilde{M} \to \widetilde{N} \to \widetilde{M}$
is the identity. Hence $\widetilde{M} \to \widetilde{N}$ is an isomorphism.
Let $K = \Ker(M \to N)$ and $Q = \Coker(M \to N)$.
Recall that the functor $M \mapsto \widetilde{M}$ is exact, see
Constructions, Lemma \ref{constructions-lemma-proj-sheaves}.
Hence we see that $\widetilde{K} = 0$ and $\widetilde{Q} = 0$.
Recall that $A$ is a Noetherian ring, $M$ is a finitely generated
$A$-module, and $N$ is a graded $A$-module such that
$N' = \bigoplus_{n \geq 0} N_n$ is finitely generated by
the last part of Lemma \ref{lemma-coherent-on-proj}.
Hence $K' = \bigoplus_{n \geq 0} K_n$ and
$Q' = \bigoplus_{n \geq 0} Q_n$ are finite $A$-modules.
Observe that $\widetilde{Q} = \widetilde{Q'}$ and
$\widetilde{K} = \widetilde{K'}$.
Thus to finish the proof it suffices to show
that a finite $A$-module $K$ with $\widetilde{K} = 0$
has only finitely many nonzero homogeneous parts $K_d$ with $d \geq 0$.
To do this, let $x_1, \ldots, x_r \in K$ be homogeneous generators
say sitting in degrees $d_1, \ldots, d_r$.
Let $f_1, \ldots, f_n \in A_1$ be elements generating $A$ over $A_0$.
For each $i$ and $j$ there exists an $n_{ij} \geq 0$ such that
$f_i^{n_{ij}} x_j = 0$ in $K_{d_j + n_{ij}}$: if not then
$x_i/f_i^{d_i} \in K_{(f_i)}$ would not be zero, i.e., $\widetilde{K}$
would not be zero.
Then we see that $K_d$ is zero for $d > \max_j(d_j + \sum_i n_{ij})$
as every element of $K_d$ is a sum of terms where each term is a
monomials in the $f_i$ times one of the $x_j$ of total degree $d$.
\end{proof}

\noindent
Let $A$ be a graded ring such that $A_0$ is Noetherian and $A$ is generated
by finitely many elements of $A_1$ over $A_0$. Recall that
$A_+ = \bigoplus_{n > 0} A_n$ is the irrelevant ideal.
Let $M$ be a graded $A$-module. Recall that
$M$ is an $A_+$-power torsion module if for all $x \in M$
there is an $n \geq 1$ such that $(A_+)^n x = 0$, see
More on Algebra, Definition \ref{more-algebra-definition-f-power-torsion}.
If $M$ is finitely generated, then we see that
this is equivalent to $M_n = 0$ for $n \gg 0$.
Sometimes $A_+$-power torsion modules are called torsion modules.
Sometimes a graded $A$-module $M$ is called torsion free if
$x \in M$ with $(A_+)^n x = 0$, $n > 0$ implies $x = 0$.
Denote $\text{Mod}_A$ the category of graded $A$-modules,
$\text{Mod}^{fg}_A$ the full subcategory of finitely generated ones,
and $\text{Mod}^{fg}_{A, torsion}$ the full subcategory of
modules $M$ such that $M_n = 0$ for $n \gg 0$.

\begin{proposition}
\label{proposition-coherent-modules-on-proj}
Let $A$ be a graded ring such that $A_0$ is Noetherian and $A$ is generated
by finitely many elements of $A_1$ over $A_0$.
Set $X = \text{Proj}(A)$. The functor $M \mapsto \widetilde M$
induces an equivalence
$$
\text{Mod}^{fg}_A/\text{Mod}^{fg}_{A, torsion}
\longrightarrow
\textit{Coh}(\mathcal{O}_X)
$$
whose quasi-inverse is given by
$\mathcal{F} \longmapsto \bigoplus_{n \geq 0} \Gamma(X, \mathcal{F}(n))$.
\end{proposition}

\begin{proof}
The subcategory $\text{Mod}^{fg}_{A, torsion}$ is a Serre subcategory
of $\text{Mod}^{fg}_A$, see
Homology, Definition \ref{homology-definition-serre-subcategory}.
This is clear from the description of objects given above but it also follows
from More on Algebra, Lemma \ref{more-algebra-lemma-I-power-torsion}.
Hence the quotient category on the left of the arrow is defined
in Homology, Lemma \ref{homology-lemma-serre-subcategory-is-kernel}.
To define the functor of the proposition, it suffices to show that
the functor $M \mapsto \widetilde M$ sends torsion modules to $0$.
This is clear because for any $f \in A_+$ homogeneous the
module $M_f$ is zero and hence the value $M_{(f)}$ of $\widetilde M$
on $D_+(f)$ is zero too.

\medskip\noindent
By Lemma \ref{lemma-coherent-on-proj} the proposed quasi-inverse
makes sense. Namely, the lemma shows that
$\mathcal{F} \longmapsto \bigoplus_{n \geq 0} \Gamma(X, \mathcal{F}(n))$
is a functor $\textit{Coh}(\mathcal{O}_X) \to \text{Mod}^{fg}_A$
which we can compose with the quotient functor
$\text{Mod}^{fg}_A \to \text{Mod}^{fg}_A/\text{Mod}^{fg}_{A, torsion}$.

\medskip\noindent
By Lemma \ref{lemma-recover-tail-graded-module}
the composite left to right to left is isomorphic to the identity functor.

\medskip\noindent
Finally, let $\mathcal{F}$ be a coherent $\mathcal{O}_X$-module.
Set $M = \bigoplus_{n \in \mathbf{Z}} \Gamma(X, \mathcal{F}(n))$
viewed as a graded $A$-module, so that our functor sends $\mathcal{F}$ to
$M_{\geq 0} = \bigoplus_{n \geq 0} M_n$.
By Properties, Lemma \ref{properties-lemma-proj-quasi-coherent}
the canonical map $\widetilde M \to \mathcal{F}$
is an isomorphism. Since the inclusion map
$M_{\geq 0} \to M$ defines an isomorphism
$\widetilde{M_{\geq 0}} \to \widetilde M$ we conclude that
the composite right to left to right is isomorphic to the identity
functor as well.
\end{proof}





